% Translation of a scalar field: the temperature field phi of a coffee cup,
% displaced by a vector a, becomes a new field phi'. The cup at position x
% carries phi; the same cup displaced to x+a carries phi'. See _template.tex /
% _preamble.tex for the shared setup; this file keeps the driver switch + drawing.
% alt: A coffee cup at position r, and the same cup displaced by a vector a to
%   a new position, with temperature fields labeled phi and phi-prime
% width: 420
\ifnum\pdfoutput>0
\documentclass[tikz,border=3pt]{standalone}         % pdflatex → PDF proof
\else
\documentclass[tikz,border=3pt,dvisvgm]{standalone} % latex   → web SVG
\fi
% Shared setup for all site figures — \input by every figures/*.tex.
% Change the palette, font size, or driver logic HERE and every figure inherits
% it on next rebuild. See src/assets/blog/README.md for the full workflow.
%
% NOTE: this file is \input AFTER \documentclass (it needs amsmath/tikz loaded),
% but the \documentclass line itself carries the driver switch and must stay in
% each figure — see _template.tex. Keep figure .tex files starting from that
% template so the driver logic is never forgotten.

\usepackage{amsmath}

% --- Palette: sentinel hexes = the LIGHT-theme values from DESIGN.md. ----------
% The PDF proof / Overleaf output therefore look like the light rendering, and
% fig.sh rewrites each hex to its themed CSS variable in the web SVG:
%   ink    #000000 -> var(--color-text-main, currentColor)   (also all black strokes)
%   accent #003366 -> var(--color-accent, #003366)
%   muted  #555555 -> var(--color-text-muted, #555555)
%   figred #b91c1c -> var(--fig-red, #b91c1c)
%   figblue   #1d4ed8 -> var(--fig-blue, #1d4ed8)
%   figorange #c2410c -> var(--fig-orange, #c2410c)
\definecolor{ink}{HTML}{000000}
\definecolor{accent}{HTML}{003366}
\definecolor{muted}{HTML}{555555}
\definecolor{figred}{HTML}{B91C1C}
\definecolor{figblue}{HTML}{1D4ED8}
\definecolor{figorange}{HTML}{C2410C}

% --- Figure body font size = site body text. ---------------------------------
% The site renders body text at 18px ≈ 13.5pt. dvisvgm emits the SVG at its
% intrinsic point size and <Figure> shows it ~1:1, so to make figure labels match
% surrounding prose we set the figure's \normalsize to 13.5/16pt. Individual
% labels can still deviate with font=\small / font=\large as needed. Computer
% Modern matches the KaTeX body math, so $dg_p(v)$ reads identically in both.
\newcommand{\figfontsize}{\fontsize{13.5}{16}\selectfont}

% Apply it as the default font for every tikzpicture, plus sane global defaults.
% Dimension figures in mm (x=1mm,y=1mm) so geometry and this pt-based text stay
% independent — scale the drawing by editing coordinates, not the type size.
\tikzset{
  every picture/.style = {font=\figfontsize, line join=round},
  edge/.style          = {line width=1pt},
  hair/.style          = {line width=1pt},
}

% Required for the Stealth arrow tip below. Without it the build dies with
% "Unknown arrow tip kind 'Stealth'" and fig.sh (set -e) aborts, silently
% leaving the previously generated SVG in place.
\usetikzlibrary{arrows.meta}

\begin{document}
\begin{tikzpicture}[x=1mm, y=1mm]
  % --- Cup geometry -----------------------------------------------------------
  % A cup is a cylinder seen slightly from above: an ellipse rim (rx x ry) and a
  % body of height \ch closed by a front-only bottom arc. Filled with figbg so it
  % is opaque: the cups are drawn BEFORE the vectors and the vector tips land in
  % their bulk (see the painter's-order note below).
  \def\rx{9}    % rim semi-axis, horizontal
  \def\ry{3}    % rim semi-axis, vertical (foreshortening -> 3/4 view)
  \def\ch{16}   % cup body height

  % \cupat{x}{y} draws a cup whose RIM CENTER is at (x,y).
  \newcommand{\cupat}[2]{%
    % body: left wall down, front bottom arc, right wall up, then the rim ellipse.
    \filldraw[edge, ink, fill=figbg]
    ({#1-\rx},{#2}) -- ({#1-\rx},{#2-\ch})
    arc[start angle=180, end angle=360, x radius=\rx, y radius=\ry]
    -- ({#1+\rx},{#2}) arc[start angle=0, end angle=360, x radius=\rx, y radius=\ry];
  }

  % \steam{x}{y} draws two wisps rising from a rim center at (x,y).
  % Each wisp is an S-curve; the pair is offset so they read as a loose plume.
  % \gap lifts the wisps clear of the rim: measured from the rim CENTER, anything
  % below \ry (=3) sits on the lip ellipse itself and reads as attached to the cup
  % rather than as steam floating above it. \gap > \ry buys visible daylight.
  \def\gap{7}
  \newcommand{\steam}[2]{%
    \foreach \dxs/\hs in {-3/1, 3/0.85}{
      \draw[hair, ink, line cap=round]
      ({#1+\dxs},{#2+\gap})
      .. controls ({#1+\dxs-2.4},{#2+\gap+4*\hs}) and ({#1+\dxs+2.4},{#2+\gap+5*\hs})
      .. ({#1+\dxs},{#2+\gap+9*\hs});
    }
  }

  % --- Layout ----------------------------------------------------------------
  % Cup rim centers, and the origin of x.
  \def\rimax{6}\def\rimay{14}    % cup 1 rim center
  \def\rimbx{56}\def\rimby{34}   % cup 2 rim center
  \def\ox{-22}\def\oy{-24}       % origin point (tail of x)

  % The vectors point into the BULK of each cup, not at the rim center (which sits
  % on the cup's top face). Drop the target \bulk mm below each rim center so the
  % head lands in the middle of the visible body. Cup body spans rim y down to
  % y-\ch (=16), so ~8mm is its mid-height.
  \def\bulk{8}
  \def\ax{\rimax}\def\ay{\rimay-\bulk}
  \def\bx{\rimbx}\def\by{\rimby-\bulk}

  % --- Cups first, vectors on top ---------------------------------------------
  % Painter's order matters: the cups carry an OPAQUE figbg fill, so anything
  % drawn before them that lands inside a rim gets painted over. The vector tips
  % are meant to sit in the BULK of each cup, so the cups must go down FIRST and
  % the arrows on top of them.
  \cupat{\rimax}{\rimay}
  \steam{\rimax}{\rimay}
  \cupat{\rimbx}{\rimby}
  \steam{\rimbx}{\rimby}

  % --- Vectors: tip-to-tail addition x + a = x' -------------------------------
  % x's head lands in the bulk of cup 1; a starts at exactly that point ({\ax},{\ay})
  % and its head lands in the bulk of cup 2. Sharing the point is what makes the
  % tip-to-tail composition read, so neither arrow is shortened at that junction.
  % `round` on the tip softens the arrowhead's corners (it otherwise inherits a
  % miter join and renders with sharp points).
  \tikzset{vec/.style = {edge, ink, line cap=round,
  -{Stealth[length=2.6mm, width=2mm, round]}}}

  % x: origin -> bulk of cup 1 (shorten < just clears the origin dot).
  % NB: `shorten` takes an ABSOLUTE length and ignores the x=1mm/y=1mm scaling,
  % so it must carry an explicit `mm` unit — a bare number means points.
  \draw[vec, shorten <=1.2mm] (\ox,\oy) -- ({\ax},{\ay});
  % a: tip of x -> bulk of cup 2.
  \draw[vec] ({\ax},{\ay}) -- ({\bx},{\by});
  % origin dot.
  \fill[ink] (\ox,\oy) circle[radius=0.9];

  % --- Vector labels ----------------------------------------------------------
  % Both x and a are diagonal segments: place each label at the segment's true
  % midpoint pushed out along the segment's own perpendicular normal, so the
  % label clears the stroke and the two read as a balanced pair (see the skill's
  % note on diagonal labels).
  % Endpoints are written out numerically here: \ay/\by expand to expressions
  % (e.g. `14-8`), so bare arithmetic like (\oy+\ay)/2 would mis-parse.
  % Symbols are set at an ABSOLUTE 17pt: the preamble's `every picture` style
  % fixes an absolute \figfontsize (13.5pt), which overrides a relative \large.
  % x: (-22,-24) -> (6,6). Midpoint (-8,-9); direction (28,30), |d| = 41.0;
  % unit normal (up-left) = (-30,28)/41.0.
  \node[text=ink, font=\fontsize{17}{20}\selectfont] at ({-8 - 5.5*30/41.0}, {-9 + 5.5*28/41.0})
  {$\boldsymbol{x}$};
  % a: (6,6) -> (56,26). Midpoint (31,16); direction (50,20), |d| = 53.9;
  % unit normal (up-left) = (-20,50)/53.9.
  \node[text=ink, font=\fontsize{17}{20}\selectfont] at ({31 - 5.5*20/53.9}, {16 + 5.5*50/53.9})
  {$\boldsymbol{a}$};

  % --- Field labels: the emphasis, mapped to figred. --------------------------
  % Placed to the outside of each cup (left of cup 1, right of cup 2) so they sit
  % clear of the steam and of the arrow between the cups.
  \node[text=figred, font=\fontsize{17}{20}\selectfont] at ({\rimax-\rx-7},{\rimay-4}) {$\phi$};
  \node[text=figred, font=\fontsize{17}{20}\selectfont] at ({\rimbx+\rx+7},{\rimby-4}) {$\phi'$};
\end{tikzpicture}
\end{document}
